\documentclass[12pt, twoside]{article}
\input{../preamble}
\usepackage{float}

\title{Physics}
\author{Chris Huson}
\date{March 2026}

\fancyhead[LE]{\thepage}
\fancyhead[RO]{\thepage}
\fancyhead[LO]{La Scuola d'Italia / Huson / Physics \\* 23 March 2026}

\begin{document}

\subsubsection*{2.5 PreQuiz Solutions: Motion, Lab Skills, and Spreadsheets}

\begin{enumerate}[itemsep=0.25cm]

\item Round each value to three significant figures.
  \begin{multicols}{2}
  \begin{enumerate}[itemsep=0.25cm]
    \item 0.004817 m $\approx$ \textbf{0.00482 m}
    \item 63.249 kg $\approx$ \textbf{63.2 kg}
    \item 7.0036 s $\approx$ \textbf{7.00 s}
    \item 21,560 N $\approx$ \textbf{21,600 N}
  \end{enumerate}
  \end{multicols}
  \vspace{0.25cm}

\item Perform each operation and give the answer in scientific notation, rounded to 3 significant figures.
  \begin{enumerate}[itemsep=0.5cm]
    \item $(3.2\times10^{4}) + (7.5\times10^{3}) = 32{,}000 + 7{,}500 = 39{,}500 = \mathbf{3.95\times10^{4}}$
    \item $(8.40\times10^{-3}) \times (2.0\times10^{4}) = 8.40 \times 2.0 \times 10^{-3+4} = 16.8 \times 10^{1} = \mathbf{1.68\times10^{2}}$
    \item $\dfrac{6.3\times10^{5}}{9.0\times10^{2}} = \dfrac{6.3}{9.0}\times10^{5-2} = 0.70\times10^{3} = \mathbf{7.00\times10^{2}}$
  \end{enumerate}
  \vspace{0.25cm}

\item Convert each value. Show one line of work using unit factors.
  \begin{enumerate}[itemsep=0.75cm]
    \item $36.0~\text{km/h} \times \dfrac{1000~\text{m}}{1~\text{km}} \times \dfrac{1~\text{h}}{3600~\text{s}}$ $= \mathbf{10.0~\textbf{m/s}}$
    \item $15.0~\text{m/s} \times \dfrac{1~\text{km}}{1000~\text{m}} \times \dfrac{3600~\text{s}}{1~\text{h}}$ $= \mathbf{54.0~\textbf{km/h}}$
  \end{enumerate}
  \vspace{0.25cm}

\item A student jogs along a straight track (forward is $+$). Start $x_0 = 5.0$ m.
  \begin{enumerate}[itemsep=0.5cm]
    \item $\Delta x = 17.0 - 5.0 = \mathbf{+12.0~\textbf{m}}$
    \item $\Delta x = 9.0 - 17.0 = \mathbf{-8.0~\textbf{m}}$
    \item Total displacement $= 9.0 - 5.0 = \mathbf{+4.0~\textbf{m}}$
    \item Total distance $= |{+12.0}| + |{-8.0}| = 12.0 + 8.0 = \mathbf{20.0~\textbf{m}}$
  \end{enumerate}

\newpage
\subsubsection*{Kinematics Equations}

\textbf{Formulas:} \quad $v = v_0 + at$ \quad\quad $d = v_0 t + \frac{1}{2}at^2$ \quad\quad $v^2 = v_0^2 + 2ad$ \quad\quad $d = \frac{1}{2}(v_0 + v)t$

\vspace{0.5cm}

\item A bicycle accelerates uniformly from rest at $1.5~\text{m/s}^2$.
  \begin{enumerate}[itemsep=0.75cm]
    \item $v = v_0 + at = 0 + (1.5)(4.0) = \mathbf{6.0~\textbf{m/s}}$
    \item $d = v_0 t + \frac{1}{2}at^2 = 0 + \frac{1}{2}(1.5)(4.0)^2 = \frac{1}{2}(1.5)(16) = \mathbf{12~\textbf{m}}$
  \end{enumerate}
  \vspace{0.5cm}

\item A ball is thrown vertically upward with an initial velocity of $12~\text{m/s}$. Use $g = 10~\text{m/s}^2$.
  \begin{enumerate}[itemsep=0.75cm]
    \item At max height, $v = 0$. \quad $0 = 12 - 10t$ \quad $\Rightarrow$ \quad $t = \mathbf{1.2~\textbf{s}}$
    \item $d = v_0 t + \frac{1}{2}at^2 = (12)(1.2) + \frac{1}{2}(-10)(1.2)^2 = 14.4 - 7.2 = \mathbf{7.2~\textbf{m}}$
  \end{enumerate}

\vspace{0.5cm}

\subsubsection*{Lab and Technology Skills --- Solutions}

\item Graphical Analysis curve fit.

The data follows $x = At^2 + Bt + C$. A quadratic fit gives approximately:
$$A \approx 0.160, \quad B \approx -0.040, \quad C \approx 0.20$$

  \begin{enumerate}[itemsep=0.5cm]
    \item[(a--b)] Students create graph and report fit parameters from Graphical Analysis.
    \item[(c)] Since $A = \frac{1}{2}a$, we get $a = 2A = 2(0.160) = \mathbf{0.32~\textbf{m/s}^2}$
    \item[(d)] File saved to Google Drive with student name.
  \end{enumerate}
  \vspace{0.25cm}

\item Spreadsheet calculation.
  \begin{enumerate}[itemsep=0.25cm]
    \item[(a--d)] Students paste screenshot, enter values, and write formula.
    \item[(e)] $F = ma = (0.390)(0.32) \approx \mathbf{0.12~\textbf{N}}$

    \textit{Note: exact answer depends on student's curve fit parameters.}
  \end{enumerate}

\end{enumerate}
\end{document}
